\documentclass[12pt]{article}

\usepackage[margin=1in]{geometry} 
\usepackage{amsmath,amsthm,amssymb,scrextend}
\usepackage{xcolor}
\usepackage{pgfplots}
\usepackage{enumitem}

\newcommand{\N}{\mathbb{N}}
\newcommand{\Z}{\mathbb{Z}}
\newcommand{\I}{\mathbb{I}}
\newcommand{\R}{\mathbb{R}}
\newcommand{\Q}{\mathbb{Q}}
\renewcommand{\qed}{\hfill$\blacksquare$}
\let\newproof\proof
\renewenvironment{proof}{\begin{addmargin}[1em]{0em}\begin{newproof}}{\end{newproof}\end{addmargin}\qed}
% \newcommand{\expl}[1]{\text{\hfill[#1]}$}

\newenvironment{theorem}[2][Theorem]{\begin{trivlist}
		\item[\hskip \labelsep {\bfseries #1}\hskip \labelsep {\bfseries #2.}]}{\end{trivlist}}
\newenvironment{lemma}[2][Lemma]{\begin{trivlist}
		\item[\hskip \labelsep {\bfseries #1}\hskip \labelsep {\bfseries #2.}]}{\end{trivlist}}
\newenvironment{problem}[2][Problem]{\begin{trivlist}
		\item[\hskip \labelsep {\bfseries #1}\hskip \labelsep {\bfseries #2.}]}{\end{trivlist}}
\newenvironment{exercise}[2][Exercise]{\begin{trivlist}
		\item[\hskip \labelsep {\bfseries #1}\hskip \labelsep {\bfseries #2.}]}{\end{trivlist}}
\newenvironment{reflection}[2][Reflection]{\begin{trivlist}
		\item[\hskip \labelsep {\bfseries #1}\hskip \labelsep {\bfseries #2.}]}{\end{trivlist}}
\newenvironment{proposition}[2][Proposition]{\begin{trivlist}
		\item[\hskip \labelsep {\bfseries #1}\hskip \labelsep {\bfseries #2.}]}{\end{trivlist}}
\newenvironment{corollary}[2][Corollary]{\begin{trivlist}
		\item[\hskip \labelsep {\bfseries #1}\hskip \labelsep {\bfseries #2.}]}{\end{trivlist}}

\setlength{\parindent}{0pt}

\begin{document}

	\textbf{(i)} When $f(x)=1$, $g(x)=e^x$, $a=0$, $b=t$:
	$$\displaystyle \left( \int_0^t e^x\,dx \right)^2\le \left( \int_0^t 1\,dx \right)\left( \int_0^t e^{2x}\,dx \right)$$
	$$\implies (e^t-1)^2\le t\cdot \frac{e^{2t}-1}{2}$$
	$$\implies (e^t-1)^2\le \frac{t}{2}(e^t-1)(e^t+1)$$

	$\because t>0\quad\therefore e^t-1>0$
	$$\frac{e^t-1}{e^t+1}\le \frac{t}{2}\\$$


	
	\textbf{(ii)} When $f(x)=x$, $g(x)=e^{-\frac{x^2}{4}}$, $a=0$, $b=1$:
	$$\displaystyle \left( \int_0^1 xe^{-\frac{x^2}{4}}\,dx \right)^2\le
	\left( \int_0^1 x^2\,dx \right)\left( \int_0^1 e^{-\frac{x^2}{2}}\,dx \right)$$

	$$ \because \int_0^1 x^2\,dx=\frac13\quad\&\&\quad
	\int_0^1 xe^{-x^2/4}\,dx=\left[-2e^{-x^2/4}\right]_0^1=2(1-e^{-1/4})$$

	$$\therefore 4(1-e^{-1/4})^2\le \frac13\int_0^1 e^{-\frac{x^2}{2}}\,dx $$
	
	$$\implies \int_0^1 e^{-x^2/2}\,dx\ge 12(1-e^{-1/4})^2$$
	


	\textbf{(iii)} When $f(x)=1$, $g(x)=\sqrt{\sin x}$, $a=0$, $b=\frac{\pi}{2}$:

	$$\displaystyle \left( \int_0^{\pi/2} \sqrt{\sin x}\,dx \right)^2\le \left( \int_0^{\pi/2} 1\,dx \right)\left( \int_0^{\pi/2} \sin x\,dx \right)$$
	
	$$\implies \left( \int_0^{\pi/2} \sqrt{\sin x}\,dx \right)^2\le \frac{\pi}{2}$$
	
	$$\because \sin x>0,\quad x \in [0,\frac{\pi}{2}]$$
	$$\therefore \int_0^{\pi/2} \sqrt{\sin x}\,dx\le \sqrt{\frac{\pi}{2}}$$


	When $f(x)=\cos x$, $g(x)=(\sin x)^\frac{1}{4}$, $a=0$, $b=\frac{\pi}{2}$:

	$$\displaystyle \left( \int_0^{\pi/2} (\sin x)^{1/4} \cos x\,dx \right)^2\le \left( \int_0^{\pi/2} \cos^2 x\,dx \right)\left( \int_0^{\pi/2} \sqrt{\sin x}\,dx \right)$$

	$$\because \int_0^{\pi/2} (\sin x)^{1/4} \cos x\,dx = \left[ \frac{4}{5}(\sin x)^{5/4} \right]_0^{\pi/2} = \frac{4}{5}$$

	$$\because \int_0^{\pi/2} \cos^2 x\,dx = \int_0^{\pi/2} \frac{1+\cos(2x)}{2}\,dx = \left[ \frac{x}{2} + \frac{\sin(2x)}{4} \right]_0^{\pi/2} = \frac{\pi}{4}$$


	$$\therefore \left( \frac{4}{5} \right)^2 \le \frac{\pi}{4} \cdot \int_0^{\pi/2} \sqrt{\sin x}\,dx$$

	$$\implies \frac{16}{25} \le \frac{\pi}{4} \cdot \int_0^{\pi/2} \sqrt{\sin x}\,dx$$

	$$\implies \frac{64}{25\pi} \le \int_0^{\pi/2} \sqrt{\sin x}\,dx$$

	$$\implies \frac{64}{25\pi} \le \int_0^{\pi/2} \sqrt{\sin x}\,dx \le \sqrt{\frac{\pi}{2}}$$


	\hrule

	\begin{center}
		This solution was submitted by \textbf{Jon}!\\[4pt]
		Published and refined by \textbf{Jon}.\\[8pt]
		\textit{Last updated: April 6, 2026}
	\end{center}
	
\end{document}
